\documentclass{CUEBthesis}
% 文章信息
\titlecn{首经贸博士学位论文LaTeX模板}
\titleen{LaTeX Template of Postgraduate Thesis of \\Capital University of Economics and Business}

\department{管理工程学院}
\minormajor{管理科学与工程}
\author{姜小喵}
\authoren{Jiang Xiaomiao}
\supervisor{猫大喵\ 教授}
\supervisoren{Prof. Mao Damiao}
\subsupervisor{}
\thesisdate{year=2023,month=5}

\newif \ifblindreview % 条件语句，是否是盲审版本
% \blindreviewtrue
\blindreviewfalse 
\begin{document}
%%%%%%%%%%%%%%%%%%%%%%%%%%%%%%%%%%%%%%%%%%%%%%%%%%
% 中文封面
% -----------------------------------------------%
\makecoverpage
\cleardoublepage

%%%%%%%%%%%%%%%%%%%%%%%%%%%%%%%%%%%%%%%%%%%%%%%%%%
% 前置部分的页眉页脚设置
% -----------------------------------------------%
%%%%%%%%%%%%%%%%%%%%%%%%%%%%%%%%%%%%%%%%%%%%%%%%%%
% 英文封面
% -----------------------------------------------%

\makecoverpageen
\cleardoublepage

%%%%%%%%%%%%%%%%%%%%%%%%%%%%%%%%%%%%%%%%%%%%%%%%%%
% 声明页
% -----------------------------------------------%
\newgeometry{top=25.4mm,bottom=25.4mm,left=31.7mm,right=31.7mm}
\announcement
\restoregeometry
\cleardoublepage

% 正文和后置部分用阿拉伯数字编连续码，前置部分用罗马数字单独编连续码（封面除外）。
% 设置声明页后面的页码

\fancyhf{} % 清空原有格式
\fancypagestyle{plain}{}
\renewcommand{\headrulewidth}{0pt}
\setcounter{page}{1} % 从1开始编号页码
\pagenumbering{Roman} % 大写罗马字母
% 设置页眉和页脚
\fancyhf[ROF]{\zihao{5}\thepage}
\fancyhf[LEF]{\zihao{5}\thepage}
% 正文以前部分无需页眉

%%%%%%%%%%%%%%%%%%%%%%%%%%%%%%%%%%%%%%%%%%%%%%%%%%
% 中文摘要
% -----------------------------------------------%
\makeatletter\@openrightfalse
% 设置中文摘要
\begin{abstractcn}
摘要是学位论文内容的简短陈述，应体现论文工作的核心思想。论文摘要应力求语言精练准确，博士学位论文的中文摘要一般约800～1200字，硕士学位论文的中文摘要一般约500字。摘要内容应涉及本项科研工作的目的和意义、研究思想和方法、研究成果和结论，博士学位论文必须突出论文的创造性成果，硕士学位论文必须突出论文的新见解。论文摘要中不要出现图片、图表、表格或其他插图材料。

关键词是为用户查找文献，从文中选取出来用来揭示全文主题内容的一组词语或术语，应尽量采用词表中的规范词（参照相应的技术术语标准）。关键词一般为3～5个，按词条的外延层次排列（外延大的排在前面）。关键词之间用分号分开，最后一个关键词后不打标点符号。

\heiti 关键词：\songti 首经贸；博士学位论文；LaTeX模板

\end{abstractcn} 
\newpage

%%%%%%%%%%%%%%%%%%%%%%%%%%%%%%%%%%%%%%%%%%%%%%%%%%
% 英文摘要
% -----------------------------------------------%
\begin{abstracten}
LaTeX Template of Postgraduate Thesis of Capital University of Economics and Business LaTeX Template of Postgraduate Thesis of Capital University of Economics and Business LaTeX Template of Postgraduate Thesis of Capital University of Economics and Business LaTeX Template of Postgraduate Thesis of Capital University of Economics and Business LaTeX Template of Postgraduate Thesis of Capital University of Economics and Business LaTeX Template of Postgraduate Thesis of Capital University of Economics and Business LaTeX Template of Postgraduate Thesis of Capital University of Economics and Business LaTeX Template of Postgraduate Thesis of Capital University of Economics and Business LaTeX Template of Postgraduate Thesis of Capital University of Economics and Business LaTeX Template of Postgraduate Thesis of Capital University of Economics and Business

为了国际交流的需要，论文须有英文摘要。英文摘要的内容及关键词应与中文摘要及关键词一致，要符合英语语法，语句通顺，文字流畅。英文摘要部分的标题为“Abstract”，用 Times New Roman体三号字，居中书写，单倍行距，段前空 24 磅，段后空 18 磅。摘要内容用小四号 Times New Roman 字体书写，两端对齐，标点符号用英文标点符号。“Key Words”需加粗，与中文摘要部分的关键词对应，每个关键词之间用分号间隔。

\textbf{Key Words:~} CUEB;~ LaTeX;~ Template

\end{abstracten} 
\newpage

%%%%%%%%%%%%%%%%%%%%%%%%%%%%%%%%%%%%%%%%%%%%%%%%%%
% 正文页眉页脚
% -----------------------------------------------%
%\setcounter{page}{1} % 重置目录页码为小写罗马字体
%\pagenumbering{roman} % 设置页码为小写罗马字体
% 设置页眉和页脚 %
\pagestyle{fancy}
\fancyhf[ROF]{\zihao{5}\thepage}
\fancyhf[LEF]{\zihao{5}\thepage}  % 所有（奇数和偶数）中间页脚

% 目录
% -------------------------------------------%
\setlength{\baselineskip}{20pt}
\tableofcontents
\cleardoublepage

% 图表目录 （没见有人写过）
%\setlength{\baselineskip}{20pt}
%{%
%\let\oldnumberline\numberline
%\renewcommand{\numberline}{\figurename~\oldnumberline}
%\listoffigures
%}
%\newpage
%
%\setlength{\baselineskip}{20pt}
%{
%\let\oldnumberline\numberline
%\renewcommand{\numberline}{\tablename~\oldnumberline}
%\listoftables
%}
%\newpage


% 重新设置正文行间距，因为前置部分设置时候行间距被改过
\renewcommand*{\baselinestretch}{1.0}   % 几倍行间距
\setlength{\baselineskip}{20pt}         % 基准行间距

% 正文

%页眉页脚
\pagestyle{fancy}
\fancypagestyle{plain}{}
\renewcommand{\chaptername}{第{\thechapter}章}
\renewcommand{\chaptermark}[1]{\markboth{\chaptername\quad #1}{}}
\renewcommand{\headrulewidth}{1pt}
\fancyhf[ROH]{\songti \zihao{-5} 首都经济贸易大学博士学位论文}
\fancyhf[LEH]{\songti \zihao{-5} \leftmark} % 偶数页显示章节名
% --------------------------------------------%
\setcounter{page}{1} % 重置页码编号
\pagenumbering{arabic} % 设置页码编号为阿拉伯数字
\fancyhf[ROF]{\zihao{5}\thepage}
\fancyhf[LEF]{\zihao{5}\thepage}

\chapter{绪论}
学位论文为了反映出作者确已掌握了坚实的基础理论和系统的专门知识，具有开阔的科学视野，对研究方案作了充分论证，因此，有关历史回顾和前人工作的综述分析，以及理论分析等，可以单独成章，用足够的文字叙述。博士学位论文绪论一般不少于 0.8万字，硕士学位论文绪论一般不少于 0.5 万字。\footnote{脚注内容采用宋体小五号字，按两端对齐格式书写，单倍行距，段前段后均空 0 磅。脚注的序号按页编排，不同页的脚注序号不需要连续。详细规定见本页脚注}
\section{研究背景及意义}
正文中某句话需要具体注释、且注释内容与正文内容关系不大时可以采用脚注方式。在正文中需要注释的句子结尾处用 xxxx 样式的数字编排序号，以“上标”字体标示在需要注释的句子末尾。在当页下部书写脚注内容。\footnote{脚注内容采用宋体小五号字，按两端对齐格式书写，单倍行距，段前段后均空 0 磅。脚注的序号按页编排，不同页的脚注序号不需要连续。详细规定见本页脚注}
\section{文献综述}
（1）标注格式：引用参考文献标注方式应全文统一，标注的格式为[序号]，放在引文或转述观点的最后一个句号之前，所引文献序号用小4号Times New Roman体、以上角标形式置于方括号中，如：“…成果\upcite{pritchard1969statistical}”。
（2）序号：标注的序号可选择如下两中方式之一：根据在正文中被首次引用出现的先后次序递增，或者按第一作者姓的英文字母或拼音字母的英文字母顺序递增。如一篇论著在论文中多处引用时，以第一次出现的序号为准；当某一论述同时引证多篇参考文献时，应采用[1，2，…]的格式，标注方括号中的序号按增序排列。
（3）当提及的参考文献为文中直接说明时，其序号应该与正文排齐，如“由文献\cite{pritchard1969statistical}可知：”。

\section{研究内容与研究思路}
\subsection{研究内容}
\subsection{研究方法}
\subsection{技术路线}
\begin{figure}[ht]
\centering
\includegraphics[width=10cm]{cueb.eps}
\caption{The procedures of generating the positive and negative relation images.}
\label{smi}
\end{figure}

\begin{figure}[ht]
\centering
\includegraphics[width=10cm]{cueb.eps}
\caption{The procedures of generating the positive and negative relation images.}
\label{smi}
\end{figure}

\begin{figure}[ht]
\centering
\includegraphics[width=10cm]{cueb.eps}
\caption{The procedures of generating the positive and negative relation images.}
\label{smi}
\end{figure}

\begin{figure}[ht]
\centering
\includegraphics[width=10cm]{cueb.eps}
\caption{The procedures of generating the positive and negative relation images.}
\label{smi}
\end{figure}

\begin{figure}[ht]
\centering
\includegraphics[width=10cm]{cueb.eps}
\caption{The procedures of generating the positive and negative relation images.}
\label{smi}
\end{figure}


\begin{figure}[ht]
\centering
\includegraphics[width=10cm]{cueb.eps}
\caption{The procedures of generating the positive and negative relation images.}
\label{smi}
\end{figure}

\begin{figure}[ht]
\centering
\includegraphics[width=10cm]{cueb.eps}
\caption{The procedures of generating the positive and negative relation images.}
\label{smi}
\end{figure}

\begin{figure}[ht]
\centering
\includegraphics[width=10cm]{cueb.eps}
\caption{The procedures of generating the positive and negative relation images.}
\label{smi}
\end{figure}

\section{论文结构安排}


\chapter{格式样例}
\section{二级标题}
\subsection{三级标题}
\subsubsection{四级标题}
首都经济贸易大学创建于1956年，是由原北京经济学院和原北京财贸学院于1995年3月合并、组建的北京市属重点大学。建校60余年来，秉承“崇德尚能，经世济民”的校训精神，首经贸人将思想和行动凝聚为“博纳敏行，知行合一”的骆驼精神，坚持“立足北京、服务首都、面向全国、走向世界”，先后为国家和北京市培养了十八万优秀人才。首经贸校友遍布海内外，成为服务国家和北京市经济社会发展的重要力量。

首经贸校本部位于丰台区花乡，以全日制本科和研究生教育为主，红庙校区位于朝阳区红庙，以留学生和非全日制教育为主。

引用文献\cite{lamport1994latex}。

插入图片，如图\ref{fig:sample_figure}所示。
\begin{figure*}[htb]
\centering
\includegraphics[width=0.5\textwidth]{cueb.png}
\caption{插图样例}
\label{fig:sample_figure}
\end{figure*}

插入表格，如表\ref{tab:sample_table}所示。
\begin{table}[htp]
\small
\centering
\setlength{\abovecaptionskip}{12pt}
\setlength{\belowcaptionskip}{6pt}
\caption{The holding gains and losses and the maximum profit of different stocks.}
\begin{tabular}{p{3cm}p{5cm}p{5cm}}\toprule[1.5pt]
Code& Holding Gains and Losses ($P_{H}$)& Maximum Profit ($P_{M}$)\\\midrule[1pt]
600809& 129430.08 & 14897932089.51 \\
300330& -4252.93 & 5503839.45 \\
002580& -26536.82 & 10147989.03 \\
603808& -3273.32 & 260181.89 \\
601318& 61696.65 & 160287819.61 \\
603123& -5574.92 & 578124.70\\
\bottomrule[1.5pt]
\end{tabular}
\label{tab:sample_table}
\end{table}

公式样例如式（\ref{equ}）所示：
\begin{equation}
sin^2x+cos^2x=1
\label{equ}
\end{equation} 

\begin{equation}
sin^2x+cos^2x=1
\label{equ}
\end{equation} 

\begin{equation}
sin^2x+cos^2x=1
\label{equ}
\end{equation} 



%%%%%%%%%%%%%%%%%%%%%%%%%%%%%%%%%%%%%%%%%%%%%%%%%%
% 后续内容，章节无编号
% --------------------------------------------%

\setcounter{chapter}{3}
\setcounter{section}{0}
\fancyhf[ROH]{\songti \zihao{-5} 首都经济贸易大学博士学位论文}
\fancyhf[LEH]{\songti \zihao{-5} 结论与展望}

\begin{conclusion}
学位论文的结论单独作为一章，但不加章号。
结论是对论文主要研究结果、论点的提炼与概括，应准确、简明，完整，有条理，使人看后就能全面了解论文的意义、目的和工作内容。主要阐述自己的创造性工作及所取得的研究成果在本学术领域中的地位、作用和意义。同时，要严格区分自己取得的成果与导师及他人的科研工作成果。
在评价自己的研究工作成果时，要实事求是，除非有足够的证据表明自己的研究是“首次”的，“领先”的，“填补空白”的，否则应避免使用这些或类似词语。
\end{conclusion}


\fancyhf[ROH]{\songti \zihao{-5} 首都经济贸易大学博士学位论文}
\fancyhf[LEH]{\songti \zihao{-5} 参考文献}

\reference{thesis-references}

\fancyhf[ROH]{\songti \zihao{-5} 首都经济贸易大学博士学位论文}
\fancyhf[LEH]{\songti \zihao{-5} 附录A}

\begin{appendix}
\chapter*{附录A}
\addcontentsline{toc}{chapter}{附录A}
\renewcommand{\thetable}{A.\arabic{table}}
\renewcommand{\thefigure}{A.\arabic{figure}}

附录图片，如图\ref{appdxfig}所示。
\begin{figure*}[htb]
\centering
\includegraphics[width=0.5\textwidth]{cueb.png}
\caption{插图样例}
\label{appdxfig}
\end{figure*}

附录表格，如表\ref{appdxtab}所示。
\begin{table}[htp]
\small
\centering
\setlength{\abovecaptionskip}{12pt}
\setlength{\belowcaptionskip}{6pt}
\caption{The holding gains and losses and the maximum profit of different stocks.}
\begin{tabular}{p{3cm}p{5cm}p{5cm}}\toprule[1.5pt]
Code& Holding Gains and Losses ($P_{H}$)& Maximum Profit ($P_{M}$)\\\midrule[1pt]
600809& 129430.08 & 14897932089.51 \\
300330& -4252.93 & 5503839.45 \\
002580& -26536.82 & 10147989.03 \\
603808& -3273.32 & 260181.89 \\
601318& 61696.65 & 160287819.61 \\
603123& -5574.92 & 578124.70\\
\bottomrule[1.5pt]
\end{tabular}
\label{appdxtab}
\end{table}

附录表格，如表\ref{appdxtab}所示。
\begin{table}[htp]
\small
\centering
\setlength{\abovecaptionskip}{12pt}
\setlength{\belowcaptionskip}{6pt}
\caption{The holding gains and losses and the maximum profit of different stocks.}
\begin{tabular}{p{3cm}p{5cm}p{5cm}}\toprule[1.5pt]
Code& Holding Gains and Losses ($P_{H}$)& Maximum Profit ($P_{M}$)\\\midrule[1pt]
600809& 129430.08 & 14897932089.51 \\
300330& -4252.93 & 5503839.45 \\
002580& -26536.82 & 10147989.03 \\
603808& -3273.32 & 260181.89 \\
601318& 61696.65 & 160287819.61 \\
603123& -5574.92 & 578124.70\\
\bottomrule[1.5pt]
\end{tabular}
\label{appdxtab}
\end{table}

附录表格，如表\ref{appdxtab}所示。
\begin{table}[htp]
\small
\centering
\setlength{\abovecaptionskip}{12pt}
\setlength{\belowcaptionskip}{6pt}
\caption{The holding gains and losses and the maximum profit of different stocks.}
\begin{tabular}{p{3cm}p{5cm}p{5cm}}\toprule[1.5pt]
Code& Holding Gains and Losses ($P_{H}$)& Maximum Profit ($P_{M}$)\\\midrule[1pt]
600809& 129430.08 & 14897932089.51 \\
300330& -4252.93 & 5503839.45 \\
002580& -26536.82 & 10147989.03 \\
603808& -3273.32 & 260181.89 \\
601318& 61696.65 & 160287819.61 \\
603123& -5574.92 & 578124.70\\
\bottomrule[1.5pt]
\end{tabular}
\label{appdxtab}
\end{table}

\end{appendix} 

\fancyhf[ROH]{\songti \zihao{-5} 首都经济贸易大学博士学位论文}
\fancyhf[LEH]{\songti \zihao{-5} 附录B}

\include{content/appdxb}

\fancyhf[ROH]{\songti \zihao{-5} 首都经济贸易大学博士学位论文}
\fancyhf[LEH]{\songti \zihao{-5} 攻读博士学位期间取得的研究成果}

\begin{additional}
在学期间发表的学术论文分以下三部分按顺序分别列出，每部分之间空1行，序号可连续排列：
1．已经刊载的学术论文（本人是第一作者，或者导师为第一作者本人是第二作者）。按照参考文献格式书写，并在其后加括号，括号内注明该文检索类型，检索号，期刊的影响因子等。如果该论文未被检索，在括号内注明期刊级别或属于何种检索源刊，例如SSCI 源刊、CSSCI源刊、EI 源刊等。
2．尚未刊载，但已经接到正式录用函的学术论文（本人为第一作者，或者导师为第一作者本人是第二作者）。按照参考文献格式书写。在每一篇文献后加括号注明已被期刊录用，并注明期刊级别或属于何种检索源刊，例如SSCI 源刊、CSSCI源刊、EI 源刊等。
3．其他学术论文。可列出除上述两种情况以外的其他学术论文，但必须是已经刊载或者收到正式录用函的论文。
除学术论文外，研究成果也可以是在学期间参加的研究项目、获得专利或获奖情况等。获得专利请注明专利名称、作者、专利号；奖项请注明获奖名称、颁奖部门、获奖时间及个人在获奖者中的名次。
\section*{学术论文}
\subsection*{已经刊载的学术论文}
\noindent
\begin{enumerate}[label={[\arabic*]}]
\item \textbf{Xiaomiao Jiang}, Director, Someone. CSU Latex Template[J]. CSU player: 1(1):1-10. （SCI 检索，JCR 1 区）
\item Director, \textbf{Xiaomiao Jiang}, Someone, Someother. XXXXXX[J]. Transactions on Image Processing.（SCI Under Review，JCR 1 区）
\item Director, \textbf{Xiaomiao Jiang}, Someone, Someother. XXXXXX[J]. Transactions on Circuits and Systems for Video Technology.（SCI Under Review，JCR 1 区）
\end{enumerate}


\subsection*{尚未刊载的学术论文}
\noindent
\begin{enumerate}[label={[\arabic*]}]
\item \textbf{Xiaomiao Jiang}, Director, Someone. CSU Latex Template[J]. CSU player: 1(1):1-10. （SCI 检索，JCR 1 区）
\item Director, \textbf{Xiaomiao Jiang}, Someone, Someother. XXXXXX[J]. Transactions on Image Processing.（SCI Under Review，JCR 1 区）
\item Director, \textbf{Xiaomiao Jiang}, Someone, Someother. XXXXXX[J]. Transactions on Circuits and Systems for Video Technology.（SCI Under Review，JCR 1 区）
\end{enumerate}

\subsection*{其他学术论文}
\noindent
\begin{enumerate}[label={[\arabic*]}]
\item \textbf{Xiaomiao Jiang}, Director, Someone. CSU Latex Template[J]. CSU player: 1(1):1-10. （SCI 检索，JCR 1 区）
\item Director, \textbf{Xiaomiao Jiang}, Someone, Someother. XXXXXX[J]. Transactions on Image Processing.（SCI Under Review，JCR 1 区）
\item Director, \textbf{Xiaomiao Jiang}, Someone, Someother. XXXXXX[J]. Transactions on Circuits and Systems for Video Technology.（SCI Under Review，JCR 1 区）
\end{enumerate}

\section*{主持和参与的研究项目}
\noindent
\begin{enumerate}[label={[\arabic*]}]
\item 主持首都经济贸易大学博士研究生学术新人项目“xxxxxxxxx”（项目号）
\item 参与国家社会科学基金项目“xxxxxxx”（项目号）
\item 参与国家自然科学基金项目“xxxxxxx”（项目号）
\end{enumerate}

\end{additional}


\fancyhf[ROH]{\songti \zihao{-5} 首都经济贸易大学博士学位论文}
\fancyhf[LEH]{\songti \zihao{-5} 致谢}

\begin{thanks}
致谢：用于评审、答辩、审议学位及提交学校存档的论文，致谢中主要感谢指导教师、在学术方面对论文的完成有直接贡献及重要帮助的团体和人士，以及感谢给予转载和引用权的资料、图片、文献、研究思想和设想的所有者等。致谢中还可以感谢提供研究经费及实验装置的基金会或企业等单位和人士。致谢辞应谦虚诚恳，实事求是，切忌浮夸与庸俗之词。
致谢内容篇幅限一页。标题和页眉内容均为“致谢”。
\end{thanks}


\fancyhf[ROH]{\songti \zihao{-5} 首都经济贸易大学博士学位论文}
\fancyhf[LEH]{\songti \zihao{-5} 作者简介}

\begin{authorintro}
博士学位论文应该提供作者简介，主要包括：姓名、性别、出生年月日、民族、出生地；简要学历、工作经历（职务）；以及攻读学位期间获得的其它奖励（除攻读学位期间取得的研究成果之外）。
\end{authorintro}


\end{document}
